% uid: ""
% school: "lycee"
% level: "3eme"
% section: ["Economie et Gestion"]
% chapter: "Statistiques"
% title: "Âge des clients d'une agence de voyages"
% difficulty: "Moyen"
% points: null
% tags: ["statistiques", "classes", "effectifs-cumules", "interpolation", "quartiles", "ecart-type"]
% provenance: "own"
% dressing: "série groupée en classes"
% rung: 3
% batch: "2026-11-20-stat-3eco-A"

% ---STATEMENT---
Une agence de voyages a relevé l'âge de ses $80$ clients du mois :

\begin{center}
\begin{tabular}{|l|c|c|c|c|c|}
\hline
Âge (ans) & $[0\,;10[$ & $[10\,;20[$ & $[20\,;30[$ & $[30\,;40[$ & $[40\,;50[$ \\ \hline
Effectif $n_i$ & $8$ & $20$ & $28$ & $16$ & $8$ \\ \hline
\end{tabular}
\end{center}

\begin{enumerate}
  \item Calculer l'âge moyen, la variance et l'écart type (on remplace chaque classe par son centre).
  \item Dresser le tableau des effectifs cumulés croissants.
  \item Déterminer par interpolation linéaire la médiane $M_e$ et les quartiles $Q_1$, $Q_3$. Calculer l'écart interquartile.
  \item Estimer le pourcentage de clients de moins de $35$ ans.
\end{enumerate}

% ---SOLUTION---
\begin{center}
\small
\begin{tabular}{|l|c|c|c|c|c|c|}
\hline
Classe & $[0;10[$ & $[10;20[$ & $[20;30[$ & $[30;40[$ & $[40;50[$ & Total \\ \hline
Centre $x_i$ & $5$ & $15$ & $25$ & $35$ & $45$ & \\ \hline
$n_i$ & $8$ & $20$ & $28$ & $16$ & $8$ & $80$ \\ \hline
$n_ix_i$ & $40$ & $300$ & $700$ & $560$ & $360$ & $1\,960$ \\ \hline
$n_ix_i^2$ & $200$ & $4\,500$ & $17\,500$ & $19\,600$ & $16\,200$ & $58\,000$ \\ \hline
ECC (borne sup.) & $8$ & $28$ & $56$ & $72$ & $80$ & \\ \hline
\end{tabular}
\end{center}
\begin{enumerate}
  \item On remplace chaque classe par son centre (calcul approché).
  $$\overline{x}=\frac{1\,960}{80}=\ib{$24{,}5$}\text{ ans}\qquad
  V=\frac{58\,000}{80}-24{,}5^2=725-600{,}25=124{,}75\qquad\sigma\approx\ib{$11{,}17$}.$$
  \item Voir la dernière ligne du tableau : $8,\ 28,\ 56,\ 72,\ 80$.
  \item \textbf{Rappel.}\par
  \begin{prop}[Interpolation linéaire]
  Si le rang $r$ cherché se trouve dans la classe $[a;b[$ et si $c_1$, $c_2$ sont les effectifs cumulés en $a$ et en $b$, alors
  $$Q=a+(b-a)\times\frac{r-c_1}{c_2-c_1}.$$
  On suppose les valeurs uniformément réparties dans la classe.
  \end{prop}
  \emph{$Q_1$} : rang $\dfrac{80}4=20$. ECC$(10)=8<20\le$ ECC$(20)=28$, donc dans $[10;20[$ :
  $$Q_1=10+10\times\frac{20-8}{28-8}=10+10\times\frac{12}{20}=\ib{$16$}.$$
  \emph{Médiane} : rang $40$, ECC$(20)=28<40\le$ ECC$(30)=56$, classe $[20;30[$ :
  $$M_e=20+10\times\frac{40-28}{56-28}=20+\frac{120}{28}=\frac{170}7\approx\ib{$24{,}29$}.$$
  \emph{$Q_3$} : rang $60$, ECC$(30)=56<60\le$ ECC$(40)=72$, classe $[30;40[$ :
  $$Q_3=30+10\times\frac{60-56}{72-56}=30+10\times\frac4{16}=\ib{$32{,}5$}.$$
  Écart interquartile : $Q_3-Q_1=32{,}5-16=16{,}5$ ans.
  \item On cherche l'effectif cumulé en $35$, milieu de $[30;40[$ : on est à mi-chemin entre ECC$(30)=56$ et ECC$(40)=72$ :
  $$56+\frac{35-30}{40-30}\times(72-56)=56+0{,}5\times16=64.$$
  $\dfrac{64}{80}=0{,}80$ : environ $\ib{$80\,\%$}$ des clients ont moins de $35$ ans.

  \textbf{Piège :} l'interpolation donne une \emph{estimation}, car on ne connaît pas la répartition exacte à l'intérieur de la classe.
\end{enumerate}

\textbf{Idée :} dans une classe on suppose les valeurs régulièrement réparties : les quartiles se lisent par interpolation sur les effectifs cumulés.
